Para la evaluación empírica de los algoritmos se diseñó un entorno controlado sobre una arquitectura x86\_64 ejecutando Linux bajo WSL2 (Ubuntu 22.04 LTS). El código fuente fue implementado en C++ (estándar C++17) y compilado mediante \texttt{g++} versión 11.4.0 aplicando el nivel máximo de optimización por compilador (\texttt{-O3}) para minimizar el impacto de llamadas a funciones y maximizar la vectorización de instrucciones. Las mediciones temporales se obtuvieron mediante el reloj de alta precisión de la biblioteca estándar (\texttt{std::chrono::high\_resolution\_clock}), mientras que el consumo de memoria máxima residente (Peak RSS) se registró mediante las llamadas del sistema POSIX (\texttt{getrusage} a través de \texttt{sys/resource.h}).

Los conjuntos de datos experimentales se generaron de forma automatizada conforme a los parámetros del enunciado:
\begin{itemize}
    \item \textbf{Ordenamiento:} Se evaluaron arreglos con tamaños $n \in \{10, 10^3, 10^5, 10^7\}$, bajo tres ordenamientos iniciales (\textit{ascendente}, \textit{descendente} y \textit{aleatorio}) y dos dominios de valores: $D_1 = \{0, \dots, 9\}$ y $D_7 = \{0, \dots, 10^7 - 1\}$. Para cada combinación se generaron 3 muestras independientes ($a, b, c$), totalizando 72 casos de prueba.
    \item \textbf{Multiplicación de Matrices:} Se evaluaron pares de matrices cuadradas $n \times n$ con dimensiones $n \in \{2^4, 2^6, 2^8, 2^{10}\} = \{16, 64, 256, 1024\}$. Se consideraron tres estructuras matriciales (\textit{densa}, \textit{diagonal} y \textit{dispersa} con un 10\% de elementos no nulos) y dos dominios numéricos: $D_0 = \{0, 1\}$ y $D_{10} = \{0, \dots, 9\}$, evaluando 3 muestras independientes por caso (72 pares de matrices).
\end{itemize}